% =============================================================================
% Chapter 11: Edge Intelligence — ML Systems Lecture Slides (Volume II)
% =============================================================================
% IMAGES NEEDED (SVGs converted to PDF):
%   - cloud-vs-edge-learning.pdf    - training-memory-amplifier.pdf
%   - adaptation-strategies.pdf     - fedavg-protocol.pdf
%   - edge-device-spectrum.pdf      - federated-systems-challenges.pdf
% =============================================================================
\documentclass[aspectratio=169, 12pt]{beamer}
\usepackage{../../assets/beamerthememlsys}

\mlsyssetup{
  volume       = {Volume II},
  chapter      = {Chapter 11},
  logo         = {../../assets/img/logo-mlsysbook.png},
  instlogo     = {../../assets/img/logo-harvard.png},
  chaptertitle = {Edge Intelligence},
}

% --- Fonts ---
\usepackage{fontspec}
\setsansfont{Helvetica Neue}[
  BoldFont={Helvetica Neue Bold},
  ItalicFont={Helvetica Neue Italic},
  BoldItalicFont={Helvetica Neue Bold Italic},
]
\setmonofont{JetBrains Mono}[Scale=0.85]

% --- Packages ---
\usepackage{booktabs}
\usepackage{amsmath}

% --- Image paths ---
\graphicspath{
  {images/}
}

% --- Chapter-specific macros ---
\newcommand{\FedAvg}{\textsc{FedAvg}}

% --- Helper: safe image include ---
\newcommand{\safeimg}[2][width=\textwidth,keepaspectratio]{%
  \IfFileExists{images/#2}{\includegraphics[#1]{#2}}{%
    \IfFileExists{#2}{\includegraphics[#1]{#2}}{%
      \fbox{\parbox[c][2.5cm][c]{0.85\linewidth}{\centering\footnotesize\textcolor{midgray}{[Missing image]}}}%
    }%
  }%
}

% --- Section count for navigation ---
\setcounter{mlsystotalsections}{8}

\title{Edge Intelligence}
\author{Vijay Janapa Reddi}
\institute{Harvard University}
\date{}

\begin{document}

% =============================================================================
% TITLE SLIDE
% =============================================================================
\mlsystitle{Edge Intelligence}{On-Device Learning and Federated Coordination}{cover_edge_intelligence.png}

% =============================================================================
% LEARNING OBJECTIVES
% =============================================================================
\begin{frame}{Learning Objectives}
\note{[2 min] Read through objectives. Emphasize the shift from datacenter
assumptions to edge constraints. Ask: ``How many of you have a phone that
predicts your next word? That model adapts on your device.''}

\footnotesize
\begin{enumerate}\setlength\itemsep{0pt}
  \item Contrast \textbf{cloud training}, \textbf{on-device learning}, and \textbf{federated learning} by data flow and privacy
  \item Quantify edge training overhead: \textbf{memory} (3--5$\times$), compute (2--3$\times$), energy (10--50$\times$)
  \item Select \textbf{adaptation strategies} (weight freezing, LoRA, sparse) matched to device constraints
  \item Apply \textbf{data efficiency techniques} (few-shot, replay) for limited local data
  \item Design \textbf{federated learning} systems with privacy-preserving aggregation across heterogeneous fleets
  \item Integrate on-device learning into production MLOps with distributed monitoring
\end{enumerate}

\end{frame}

% =============================================================================
\section{Why Edge Learning?}
% =============================================================================

\begin{frame}{The Edge Learning Paradigm}
\note{[3 min] Open with the voice assistant example. A model trained on
millions of generic samples cannot understand a regional accent. The key
insight: centralized training cannot anticipate every local context.
Ask: ``What happens to your keyboard predictions if you switch languages?''}

\footnotesize
\begin{columns}[T]
  \begin{column}{0.55\textwidth}
    A voice assistant trained on generic audio struggles with:
    \begin{itemize}\setlength\itemsep{0pt}
      \item Regional accents and vocabulary
      \item Household-specific names and patterns
      \item Evolving user behavior over time
    \end{itemize}

    \vspace{0.1cm}
    \begin{mlsyscard}{crimson}
    {\scriptsize\textbf{Core tension:} Centralized training cannot anticipate every deployment context. Models must adapt \emph{where they operate}.}
    \end{mlsyscard}

    \vspace{0.05cm}
    {\scriptsize Four drivers: \textbf{personalization}, \textbf{latency}, \textbf{privacy}, \textbf{efficiency}.}
  \end{column}
  \begin{column}{0.42\textwidth}
    \safeimg[width=\textwidth,height=4.2cm,keepaspectratio]{cloud-vs-edge-learning.pdf}
  \end{column}
\end{columns}

\end{frame}

\begin{frame}{When On-Device Learning is Necessary}
\note{[2 min] Not every application needs on-device learning. Walk through
the latency budgets that make cloud round-trips infeasible.
Common error: teams over-engineer when simpler alternatives suffice.}

\footnotesize
\begin{columns}[T]
  \begin{column}{0.48\textwidth}
    \textbf{Simpler alternatives often suffice:}
    \begin{itemize}\setlength\itemsep{0pt}
      \item Feature-based personalization (stored prefs)
      \item Cloud fine-tuning with differential privacy
      \item User-specific lookup tables + global model
    \end{itemize}

    \vspace{0.15cm}
    \mlsysalert{Decision rule:}{On-device learning only when simpler alternatives fail the requirements.}
  \end{column}
  \begin{column}{0.48\textwidth}
    \textbf{Latency makes cloud infeasible:}

    \vspace{0.1cm}
    \renewcommand{\arraystretch}{1.15}
    {\scriptsize
    \begin{tabular}{@{}lr@{}}
      \toprule
      \textbf{Application} & \textbf{Budget} \\
      \midrule
      Camera processing  & $<$33 ms \\
      Voice response     & $<$500 ms \\
      AR/VR motion-to-photon & $<$20 ms \\
      Safety-critical control & $<$10 ms \\
      \midrule
      \textcolor{errorstroke}{Network round-trip} & \textcolor{errorstroke}{50--200 ms} \\
      \bottomrule
    \end{tabular}
    }

    \vspace{0.1cm}
    \mlsysconcept{Physics decides:}{When the latency budget $<$ round-trip time, the model \emph{must} run locally.}
  \end{column}
\end{columns}

\end{frame}

\begin{frame}{Real-World Applications}
\note{[2 min] Survey the application domains. Each domain stresses a
different constraint. Gboard: first federated learning at billion-device
scale, 10--20\% accuracy improvement over static models.}

\scriptsize
\renewcommand{\arraystretch}{1.1}
\begin{tabular}{@{}lp{3cm}p{2.5cm}p{3cm}@{}}
  \toprule
  \textbf{Domain} & \textbf{Example} & \textbf{Driver} & \textbf{Key Constraint} \\
  \midrule
  \textbf{Mobile Input} & Gboard (1B+ devices) & Personalization & 200--300 MB memory \\
  \textbf{Wearables} & Health monitors & Privacy (HIPAA) & Sub-watt power \\
  \textbf{Voice} & Wake-word detect & Latency ($<$100 ms) & 1 mW always-on \\
  \textbf{Industrial IoT} & Pipeline monitoring & Connectivity & Offline days/weeks \\
  \textbf{Vision} & AR/VR, robotics & Latency ($<$20 ms) & Real-time frames \\
  \bottomrule
\end{tabular}

\vspace{0.1cm}
\begin{mlsyscard}{datastroke}
{\scriptsize\textbf{Common pattern:} Deployment variation that centralized training cannot anticipate. Learning under severe resource constraints.}
\end{mlsyscard}

\end{frame}

% =============================================================================
\section{Design Constraints}
% =============================================================================

\begin{frame}{Training is Not Inference}
\note{[3 min] The most important quantitative slide. Training amplifies
every constraint. Walk through the table column by column.
Key: memory goes up 3--5$\times$, energy goes up 10--50$\times$.
Ask: ``Why is energy amplification so much worse than memory?''}

% --- Layout: FULL-WIDTH table ---
\small
\textbf{On-device training amplifies inference constraints multiplicatively:}

\vspace{0.15cm}
\renewcommand{\arraystretch}{1.15}
{\scriptsize
\begin{tabular}{@{}lp{2.8cm}p{3.5cm}p{2.8cm}@{}}
  \toprule
  \textbf{Dimension} & \textbf{Inference} & \textbf{Training} & \textbf{Amplification} \\
  \midrule
  \textbf{Memory} & Weights + 1 activation map & Weights + activations + gradients + optimizer & \textcolor{errorstroke}{\textbf{3--5$\times$}} \\
  \textbf{Compute} & Forward pass only & Forward + backward + update & \textcolor{errorstroke}{\textbf{2--3$\times$}} \\
  \textbf{Bandwidth} & Sequential weight reads & Bidirectional gradient flow & \textcolor{errorstroke}{\textbf{5--10$\times$}} \\
  \textbf{Energy} & Single inference & Multiple gradient steps & \textcolor{errorstroke}{\textbf{10--50$\times$}} \\
  \bottomrule
\end{tabular}
}

\vspace{0.15cm}
\mlsysalert{Consequence:}{A smartphone keyboard uses 25\% of available memory for a \emph{single} training step. Compression is not optional --- it is a survival requirement.}

\end{frame}

\begin{frame}{The Training Memory Amplifier}
\note{[2 min] Show the bar chart. Emphasize the stacked components:
activations dominate for larger models. At 100M params, training needs
3.6 GB vs 400 MB for inference --- a 9$\times$ amplification.
If short: focus on the 100M example only.}

% --- Layout: FULL-WIDTH image ---
\centering
\safeimg[width=0.88\textwidth,height=5.2cm,keepaspectratio]{training-memory-amplifier.pdf}

\vspace{0.1cm}
{\small At 100M parameters: inference = 400 MB, training = \textbf{3.6 GB}. Activations dominate.}

\end{frame}

% --- ACTIVE LEARNING 1: Predict ---
\begin{frame}{Predict: How Would You Adapt a Model on a Smartwatch?}
\note{[2 min] Prediction exercise. Give students 60 seconds. Key insight:
you cannot update all parameters. The answer will be weight freezing or
bias-only adaptation. Do NOT reveal yet --- primes the adaptation section.}

\centering
\vspace{1.0cm}
{\Large\bfseries Think--Write--Share}

\vspace{0.5cm}
{\large A 10M-parameter health model must adapt on a smartwatch\\
with 500 MB total system memory and 0.5 W power budget.}

\vspace{0.5cm}
{\normalsize What would you freeze? What would you update?}

\vspace{0.5cm}
{\small\textcolor{midgray}{Write your strategy in 2 sentences. (60 seconds)}}

\end{frame}

% =============================================================================
\section{Device Spectrum}
% =============================================================================

\begin{frame}{The Edge Device Spectrum}
\note{[3 min] Walk through the four tiers. Key insight: 6+ orders of
magnitude separate the tiers. Each requires fundamentally different
adaptation strategies. Connect to Vol1 deployment spectrum.
Ask: ``Can you use the same LoRA adapter on a smartwatch and a phone?''}

% --- Layout: FULL-WIDTH image ---
\centering
\safeimg[width=0.92\textwidth,height=5.5cm,keepaspectratio]{edge-device-spectrum.pdf}

\vspace{0.1cm}
{\small Each tier demands a different adaptation strategy matched to its resource envelope.}

\end{frame}

\begin{frame}{NPU: The Silicon Dividend}
\note{[2 min] Mobile NPUs provide 20$\times$ speedup and 50$\times$ energy
savings over CPU execution. Without the NPU, on-device training drains
the battery in minutes. The NPU freed thermal headroom for backpropagation.}

\small
\begin{columns}[T]
  \begin{column}{0.55\textwidth}
    \textbf{NPU vs.\ CPU for MobileNet inference:}

    \vspace{0.15cm}
    \renewcommand{\arraystretch}{1.2}
    {\footnotesize
    \begin{tabular}{@{}lrr@{}}
      \toprule
      & \textbf{CPU} & \textbf{NPU} \\
      \midrule
      Latency      & 400 ms & 20 ms \\
      Energy       & 0.8 J  & 0.016 J \\
      \midrule
      \textbf{Speedup}    & ---   & \textbf{20$\times$} \\
      \textbf{Energy gain} & ---  & \textbf{50$\times$} \\
      \bottomrule
    \end{tabular}
    }

    \vspace{0.15cm}
    \mlsysinsight{Insight:}{The NPU does not just make inference faster --- it makes on-device training \emph{feasible}. The 50$\times$ energy gain provides thermal headroom for backpropagation bursts.}
  \end{column}
  \begin{column}{0.42\textwidth}
    \begin{mlsyscard}{computestroke}
    \textbf{The Mobile Memory Wall}\\[0.1cm]
    {\footnotesize While NPUs provide raw compute (TOPS), on-device models are \textbf{bandwidth-bound}. The 30--50$\times$ gap between mobile LPDDR5 and datacenter HBM3 makes quantization a mandatory survival strategy.}
    \end{mlsyscard}

    \vspace{0.1cm}
    {\footnotesize
    \renewcommand{\arraystretch}{1.1}
    \begin{tabular}{@{}lr@{}}
      Mobile LPDDR5 & 51 GB/s \\
      H100 HBM3    & 3,350 GB/s \\
      \textcolor{errorstroke}{\textbf{Gap}} & \textcolor{errorstroke}{\textbf{65$\times$}} \\
    \end{tabular}
    }
  \end{column}
\end{columns}

\end{frame}

% =============================================================================
\section{Model Adaptation}
% =============================================================================

\mlsysfocus{The Adaptation Principle}{%
Freeze what works globally.\\[0.3cm]
Adapt only what matters locally.%
}

\begin{frame}{Three Adaptation Strategies}
\note{[3 min] Walk through the progression from lightweight to expressive.
Each strategy trades expressiveness for resource cost.
Key: bias-only for microcontrollers, LoRA for smartphones, sparse for edge servers.}

% --- Layout: FULL-WIDTH image ---
\centering
\safeimg[width=0.92\textwidth,height=5.2cm,keepaspectratio]{adaptation-strategies.pdf}

\vspace{0.1cm}
{\small Strategy selection follows device constraints: more resources enable more expressive adaptation.}

\end{frame}

\begin{frame}{Weight Freezing: Bias-Only Adaptation}
\note{[2 min] The simplest strategy. Freeze all weights, update only biases.
Reduces trainable parameters by 100--1000$\times$. Show the equation:
only the bias gradient is non-zero. Works on microcontrollers.}

\footnotesize
\begin{columns}[T]
  \begin{column}{0.50\textwidth}
    \textbf{Standard layer:} $y = Wx + b$

    \vspace{0.1cm}
    In bias-only adaptation:
    $$\frac{\partial \mathcal{L}}{\partial W} = 0, \quad \frac{\partial \mathcal{L}}{\partial b} \neq 0$$

    \vspace{-0.1cm}
    \begin{mlsyscard}{datastroke}
    {\scriptsize \textbf{100--1000$\times$} parameter reduction. Enables training on devices without FP units.}
    \end{mlsyscard}
  \end{column}
  \begin{column}{0.47\textwidth}
    \begin{mlsyscard}{computestroke}
    {\scriptsize\texttt{%
    \textcolor{midgray}{\# Freeze all parameters}\\
    for p in model.parameters():\\
    \hspace{1em}p.requires\_grad = False\\[0.1cm]
    \textcolor{midgray}{\# Enable bias gradients only}\\
    for name, p in model.params():\\
    \hspace{1em}if "bias" in name:\\
    \hspace{2em}p.requires\_grad = True%
    }}
    \end{mlsyscard}
    {\scriptsize\textcolor{midgray}{PyTorch: 4 lines for 1000$\times$ param reduction.}}
  \end{column}
\end{columns}

\end{frame}

\begin{frame}{LoRA: Low-Rank Adaptation}
\note{[3 min] LoRA decomposes weight updates as a low-rank product.
Instead of updating the full $d \times d$ matrix, we learn two small
matrices $A$ and $B$ with rank $r \ll d$. Storage savings: 200$\times$
per user context. This enables 10 user contexts in 42 MB instead of 400 MB.}

\footnotesize
\begin{columns}[T]
  \begin{column}{0.52\textwidth}
    \textbf{Key idea:} Decompose weight updates as low-rank product.

    \vspace{0.05cm}
    $$W' = W_0 + \Delta W = W_0 + AB$$
    {\scriptsize where $A \in \mathbb{R}^{d \times r}$, $B \in \mathbb{R}^{r \times d}$, $r \ll d$}

    \vspace{0.1cm}
    \renewcommand{\arraystretch}{1.05}
    {\scriptsize
    \begin{tabular}{@{}lr@{}}
      \toprule
      \textbf{Metric} & \textbf{Value} \\
      \midrule
      Rank ($r$)         & 4--64 \\
      Trainable params   & 1--5\% \\
      Per-context saving & 200$\times$ \\
      \bottomrule
    \end{tabular}
    }

    \vspace{0.05cm}
    \mlsysconcept{Storage:}{10 contexts = 42 MB (vs.\ 400 MB full).}
  \end{column}
  \begin{column}{0.45\textwidth}
    \begin{mlsyscard}{routingstroke}
    {\scriptsize\textbf{Why LoRA works:} Weight updates have low intrinsic rank. Rank-4 captures most adaptation signal with orders-of-magnitude fewer parameters.}
    \end{mlsyscard}

    \vspace{0.05cm}
    \begin{mlsyscard}{errorstroke}
    {\scriptsize\textbf{Trade-off:} Low rank limits expressiveness. Complex domain shifts may need higher rank or sparse updates.}
    \end{mlsyscard}
  \end{column}
\end{columns}

\end{frame}

% --- ACTIVE LEARNING 2: Exercise ---
\begin{frame}{Your Turn: Personalization Storage Budget}
\note{[3 min] Give students 90 seconds. Answer: Full = 10$\times$40 = 400 MB.
LoRA = 40 + 10$\times$0.2 = 42 MB. Savings = 400/42 = 9.5$\times$.
Per-context = 40/0.2 = 200$\times$.}

\small
\begin{columns}[T]
  \begin{column}{0.55\textwidth}
    {\normalsize\bfseries Storage Budget for 10 User Contexts}

    \vspace{0.2cm}
    A 10M parameter vision model (40 MB in FP32).\\
    You need 10 user contexts on one phone.

    \vspace{0.15cm}
    \begin{itemize}
      \item Full fine-tuning: 10 copies $\times$ 40 MB = ?
      \item LoRA (rank 64, 50K adapter params): backbone + 10 adapters = ?
    \end{itemize}

    \vspace{0.1cm}
    \textbf{Calculate} total storage and per-context savings.

    {\footnotesize\textcolor{midgray}{(90 seconds --- compare with a neighbor)}}
  \end{column}
  \begin{column}{0.42\textwidth}
    \pause
    \begin{mlsyscard}{datastroke}
    \textbf{Solution:}\\[0.1cm]
    {\footnotesize
    Full: 10 $\times$ 40 = \textbf{400 MB}\\
    LoRA: 40 + 10 $\times$ 0.2 = \textbf{42 MB}\\[0.1cm]
    Total savings: \textbf{9.5$\times$}\\
    Per-context: \textbf{200$\times$}\\[0.1cm]
    \textcolor{datastroke}{\textbf{Modularity is the only way}} to scale personalization on edge devices.
    }
    \end{mlsyscard}
  \end{column}
\end{columns}

\end{frame}

% =============================================================================
\section{Data Efficiency}
% =============================================================================

\begin{frame}{Learning from Limited Local Data}
\note{[3 min] Edge devices have small, noisy, non-IID datasets.
Three strategies: few-shot learning, experience replay, data compression.
Key: prevent catastrophic forgetting while adapting to new patterns.}

\small
\begin{columns}[T]
  \begin{column}{0.48\textwidth}
    \textbf{Edge data is unlike cloud data:}
    \begin{itemize}\setlength\itemsep{0pt}
      \item Small: dozens to hundreds of samples
      \item Non-IID: reflects one user's patterns
      \item Noisy: real-world sensor data
      \item Streaming: arrives continuously
    \end{itemize}

    \vspace{0.1cm}
    \mlsysalert{Risk:}{\textbf{Catastrophic forgetting} --- new data erases old knowledge.}
  \end{column}
  \begin{column}{0.48\textwidth}
    \renewcommand{\arraystretch}{1.05}
    {\scriptsize
    \begin{tabular}{@{}lp{3.2cm}@{}}
      \toprule
      \textbf{Strategy} & \textbf{How It Helps} \\
      \midrule
      \textbf{Few-shot} & 5--20 examples via meta-learning \\
      \textbf{Replay} & Mix old + new examples in training \\
      \textbf{Contrastive} & Representations without labels \\
      \textbf{Compression} & Maximize information per bit \\
      \bottomrule
    \end{tabular}
    }

    \vspace{0.05cm}
    \mlsysinsight{Key:}{Combine replay with freezing to preserve global knowledge.}
  \end{column}
\end{columns}

\end{frame}

% =============================================================================
\section{Federated Learning}
% =============================================================================

\begin{frame}{The Federated Learning Idea}
\note{[3 min] Motivate with the hospital example: multiple hospitals want
to train an AI on patient records but cannot share raw data. Federated
learning moves the compute to the data. Only model updates are shared.
78$\times$ bandwidth savings vs. raw data upload.}

\footnotesize
\begin{columns}[T]
  \begin{column}{0.55\textwidth}
    \textbf{Problem:} 10M devices each with unique data.\\
    \textbf{Constraint:} Raw data cannot leave the device.

    \vspace{0.1cm}
    \begin{mlsyscard}{crimson}
    {\scriptsize\textbf{Federated Learning:} Move \emph{compute} to the data. Devices train locally, share only model updates.}
    \end{mlsyscard}

    \vspace{0.05cm}
    {\scriptsize\textbf{Bandwidth savings:}\\
    Raw data: 195 MB/week | Update: 2.5 MB\\
    \textcolor{datastroke}{\textbf{78$\times$ reduction}}\\[0.05cm]
    Gboard: 1B+ devices, 10--20\% accuracy gain.}
  \end{column}
  \begin{column}{0.42\textwidth}
    \safeimg[width=\textwidth,height=4.8cm,keepaspectratio]{fedavg-protocol.pdf}
  \end{column}
\end{columns}
\mlsyscite{McMahan et al., ``Communication-Efficient Learning,'' AISTATS 2017}

\end{frame}

\begin{frame}{FedAvg: The Core Protocol}
\note{[3 min] Walk through the four phases. Key equation: weighted average
by local sample count. Emphasize: raw data never leaves the device.
Common error: students think ``averaging weights'' is trivial --- it is
not when data is non-IID.}

\footnotesize
\textbf{Four phases per round:}

\vspace{0.05cm}
\renewcommand{\arraystretch}{1.05}
{\scriptsize
\begin{tabular}{@{}clp{6cm}@{}}
  \toprule
  & \textbf{Step} & \textbf{Action} \\
  \midrule
  1 & \textcolor{computestroke}{\textbf{Broadcast}} & Server sends global model $\theta^t$ to selected devices \\
  2 & \textcolor{datastroke}{\textbf{Local Train}} & Each device runs $E$ epochs of SGD on local data \\
  3 & \textcolor{routingstroke}{\textbf{Upload}} & Devices send updates $\theta_k^t$ (not raw data) \\
  4 & \textcolor{crimson}{\textbf{Aggregate}} & Server computes weighted average \\
  \bottomrule
\end{tabular}
}

\vspace{0.1cm}
\centering
$$\theta^{t+1} = \sum_{k \in \mathcal{K}_t} \frac{n_k}{n_{\mathcal{K}_t}} \theta_k^t$$

\vspace{0.05cm}
{\footnotesize Weighted by local sample count $n_k$. Privacy: \textbf{raw data never leaves the device.}}

\end{frame}

% --- ACTIVE LEARNING 3: Discussion ---
\begin{frame}{Discussion: Is FedAvg Really Private?}
\note{[3 min] Turn-and-talk. Key insight: NO --- gradient updates can be
inverted to reconstruct user data. Differential privacy (epsilon 1--8)
and secure aggregation are required. The 1--5\% accuracy cost is the
``privacy tax.'' Cold-call 2--3 pairs.}

\centering
\vspace{0.8cm}
{\large\bfseries Turn and Talk \textcolor{midgray}{(90 seconds)}}

\vspace{0.5cm}
{\large Federated learning keeps raw data on the device.\\[0.2cm]
\alert{Does this automatically guarantee privacy?}}

\vspace{0.5cm}
{\normalsize Consider: What can an adversary learn from\\
the model \emph{updates} alone?}

\vspace{0.5cm}
{\small\textcolor{lightgray}{Hint: gradient inversion attacks can reconstruct training images.}}

\end{frame}

\begin{frame}{Federated Systems Challenges at Scale}
\note{[3 min] The algorithmic elegance of FedAvg masks substantial
systems challenges. Walk through all four: stragglers (20--40\% offline),
non-IID data (10--50\% accuracy drop), communication constraints,
and privacy requirements. Each has protocol-level mitigations.}

% --- Layout: FULL-WIDTH image ---
\centering
\safeimg[width=0.92\textwidth,height=5.5cm,keepaspectratio]{federated-systems-challenges.pdf}

\vspace{0.1cm}
{\small Federated learning at scale is a \textbf{systems engineering problem}, not just an algorithm.}

\end{frame}

\begin{frame}{Communication-Efficient Updates}
\note{[2 min] Bandwidth is the dominant bottleneck. Three techniques:
gradient quantization (8--32$\times$), sparsification (10--100$\times$),
split architectures. LTE upload is 5--10 Mbps --- a 50 MB update takes
40--80 seconds and 2--3\% of battery.}

\footnotesize
\begin{columns}[T]
  \begin{column}{0.48\textwidth}
    \textbf{Wireless upload asymmetry:}\\
    {\scriptsize LTE: 5--10 Mbps up vs.\ 50+ Mbps down}

    \vspace{0.1cm}
    \renewcommand{\arraystretch}{1.05}
    {\scriptsize
    \begin{tabular}{@{}lrr@{}}
      \toprule
      \textbf{Technique} & \textbf{Compress} & \textbf{Acc.\ loss} \\
      \midrule
      Quantize (INT8)  & 8--16$\times$ & $<$1\% \\
      SignSGD (1-bit)   & 32$\times$    & 1--3\% \\
      Top-$k$ sparse    & 10--100$\times$ & $<$2\% \\
      \bottomrule
    \end{tabular}
    }

    \vspace{0.05cm}
    {\scriptsize 50 MB update over LTE: 40--80s, 2--3\% battery.}
  \end{column}
  \begin{column}{0.48\textwidth}
    \begin{mlsyscard}{computestroke}
    {\scriptsize\textbf{Split architecture:} Shared global backbone + private local head. Sync only the shared part.}
    \end{mlsyscard}

    \vspace{0.05cm}
    \begin{mlsyscard}{errorstroke}
    {\scriptsize\textbf{LoRaWAN:} 50 kbps, 1\% duty cycle. Uncompressed updates are \emph{physically impossible}. Compression is mandatory.}
    \end{mlsyscard}
  \end{column}
\end{columns}

\end{frame}

% =============================================================================
\section{Production Integration}
% =============================================================================

\begin{frame}{MLOps for On-Device Learning}
\note{[3 min] Production integration is where theory meets reality.
Key differences from centralized MLOps: tiered deployment, multi-dimensional
versioning, distributed monitoring, and resource-aware scheduling.
20--40\% of devices offline at any time.}

\footnotesize
\begin{columns}[T]
  \begin{column}{0.48\textwidth}
    \textbf{Pipeline transforms:}
    \begin{itemize}\setlength\itemsep{0pt}
      \item {\scriptsize Artifact: static model $\to$ policy + weights + configs}
      \item {\scriptsize Versioning: linear $\to$ hierarchical (base + adapters)}
      \item {\scriptsize Monitoring: centralized $\to$ federated analytics}
    \end{itemize}

    \vspace{0.05cm}
    \textbf{Tiered strategy:}\\
    {\scriptsize
    MCU $\to$ bias-only | Phone $\to$ LoRA | Edge $\to$ sparse}
  \end{column}
  \begin{column}{0.48\textwidth}
    \textbf{Orchestration:}
    \begin{itemize}\setlength\itemsep{0pt}
      \item {\scriptsize 20--40\% devices offline per round}
      \item {\scriptsize Train only when idle + charging + WiFi}
      \item {\scriptsize Thermal: suspend if device overheats}
      \item {\scriptsize Battery: $<$5\% per day for training}
    \end{itemize}

    \vspace{0.05cm}
    \begin{mlsyscard}{routingstroke}
    {\scriptsize\textbf{Validation:} Shadow model (frozen baseline vs.\ adapted) detects degradation.}
    \end{mlsyscard}
  \end{column}
\end{columns}

\end{frame}

% =============================================================================
\section{Wrap-Up}
% =============================================================================

\begin{frame}{Fallacies}
\note{[2 min] Three common misconceptions with quantitative evidence.}

\small
\textbf{Fallacy:} \textit{On-device learning matches cloud-based training capabilities.}\\
{\footnotesize A smartphone with 8 GB RAM and 5 W cannot replicate datacenter systems with TB of memory and kW of power. Training amplifies resource needs 3--10$\times$ over inference.}

\vspace{0.15cm}
\textbf{Fallacy:} \textit{Federated learning automatically preserves privacy.}\\
{\footnotesize Gradient inversion attacks can reconstruct training images from model updates. True privacy requires differential privacy ($\varepsilon = 1$--$8$, 1--5\% accuracy cost) and secure aggregation.}

\vspace{0.15cm}
\textbf{Fallacy:} \textit{Local adaptation always improves over generic models.}\\
{\footnotesize Adaptation with insufficient or noisy data can \emph{degrade} performance. Systems need fallback to known-good baselines when local data is inadequate.}

\end{frame}

\begin{frame}{Pitfalls}
\note{[2 min] Two operational pitfalls. Heterogeneity and orchestration
are the most underestimated challenges.}

\footnotesize
\textbf{Pitfall:} \textit{Assuming uniform hardware across the deployment fleet.}\\
Edge capabilities span 6+ orders of magnitude: 32 KB to 16 GB RAM, 48 MHz to 3 GHz, 10 $\mu$W to 5 W. A learning algorithm that works on flagship phones may fail catastrophically on IoT devices. Federated systems must handle \textbf{10,000$\times$} performance differences within a single deployment.

\vspace{0.15cm}
\textbf{Pitfall:} \textit{Underestimating orchestration complexity at scale.}\\
Federated systems must handle intermittent connectivity, version skew across thousands of model variants, participation bias (top 10\% of devices contribute to 50\% of rounds), and devices reconnecting after weeks offline. This is not a scaled-down cloud --- it is a fundamentally different engineering discipline.

\end{frame}

% --- RETRIEVAL PRACTICE ---
\begin{frame}{What Were the Key Ideas?}
\note{[2 min] Retrieval practice. Students write 90 seconds, no notes.
Walk around the room.}

\centering
\vspace{1.5cm}
{\Large\bfseries Close your notes.}

\vspace{0.8cm}
{\large Write down the \textbf{4 most important concepts} from today.}

\vspace{0.8cm}
{\normalsize\textcolor{midgray}{90 seconds --- no peeking.}}

\end{frame}

\begin{frame}{Key Takeaways}
\note{[2 min] Reveal. Walk through each bullet. Emphasize: edge is NOT
a scaled-down cloud. Quantitative anchors: 3--10$\times$ training
amplification, 78$\times$ bandwidth savings, 200$\times$ per-context
storage savings, 20--40\% offline rate.}

\footnotesize
\begin{itemize}\setlength\itemsep{1pt}
  \item \textbf{Mobile Memory Wall}: NPUs provide TOPS, but on-device models are bandwidth-bound. The 30--50$\times$ gap between mobile RAM and HBM makes quantization mandatory.
  \item \textbf{Training $\neq$ Inference}: On-device learning amplifies resource needs 3--10$\times$ due to activations, gradients, and optimizer state.
  \item \textbf{Three pillars}: Model Adaptation (bias/LoRA/sparse), Data Efficiency (few-shot/replay), Federated Coordination (privacy-preserving aggregation).
  \item \textbf{Heterogeneity is the default}: Edge spans 6+ orders of magnitude. Federated systems must handle stragglers, participation bias, and 10,000$\times$ capability differences.
  \item \textbf{Federated learning is a systems problem}: Stragglers, non-IID data, communication constraints, and privacy attacks require protocol-level mitigations.
  \item \textbf{Edge $\neq$ scaled-down cloud}: Thermal throttling, intermittent connectivity, and battery budgets make edge a fundamentally distinct discipline.
\end{itemize}

\end{frame}

\begin{frame}{References}
\note{[1 min] Point students to canonical papers.}

\small
\mlsysref{McMahan+17}{McMahan et al. ``Communication-Efficient Learning of Deep Networks from Decentralized Data.'' AISTATS 2017.}
\mlsysref{Li+20}{Li et al. ``Federated Learning: Challenges, Methods, and Future Directions.'' IEEE SP Magazine, 2020.}
\mlsysref{Bonawitz+19}{Bonawitz et al. ``Towards Federated Learning at Scale: A System Design.'' MLSys 2019.}
\mlsysref{Hu+22}{Hu et al. ``LoRA: Low-Rank Adaptation of Large Language Models.'' ICLR 2022.}
\mlsysref{Hard+18}{Hard et al. ``Federated Learning for Mobile Keyboard Prediction.'' 2018.}
\mlsysref{Cai+20}{Cai et al. ``TinyTL: Reduce Memory, Not Parameters for Efficient On-Device Learning.'' NeurIPS 2020.}

\end{frame}

\begin{frame}{Next Lecture: ML Ops at Scale}
\note{[1 min] Forward hook. We pushed intelligence to its physical limits.
Now: how do we sustain these deployments reliably? Platform engineering,
monitoring, and lifecycle management for the entire fleet.}

\footnotesize
\begin{columns}[c]
  \begin{column}{0.30\textwidth}
    \centering
    \begin{mlsyscard}{computestroke}
    {\large\bfseries Platform\\Engineering}\\[0.1cm]
    {\small CI/CD for ML fleets}
    \end{mlsyscard}
  \end{column}
  \begin{column}{0.30\textwidth}
    \centering
    \begin{mlsyscard}{routingstroke}
    {\large\bfseries Monitoring\\at Scale}\\[0.1cm]
    {\small Observability across tiers}
    \end{mlsyscard}
  \end{column}
  \begin{column}{0.30\textwidth}
    \centering
    \begin{mlsyscard}{datastroke}
    {\large\bfseries Lifecycle\\Management}\\[0.1cm]
    {\small From deploy to retire}
    \end{mlsyscard}
  \end{column}
\end{columns}

\vspace{0.3cm}
\centering
{\normalsize Deploying models is the beginning.\\
\textbf{Sustaining them at production scale is the real engineering.}}

\end{frame}

\end{document}
